\section{Related Work} \label{sec:related_work}
Our work closely relates to two lines of research: pruning the search space for CO problems and automated feature engineering with LLMs. \textcolor{black}{We emphasize that our approach is orthogonal to methods that directly generate algorithms using LLMs, which often exhibit degraded performance as graph size increases \citep{yang2023large,liu2024evolution}. Instead, we leverage the reasoning capabilities of LLMs to generate features for a learned algorithm for pruning.}

\textbf{Automated Feature Engineering via LLMs.} Automated feature engineering generates useful features from raw data without human effort. Prior work with LLMs has focused primarily on tabular data, where new features are derived from existing ones. These approaches either rely on manually designed feature-generation rules \citep{zhang2023openfe,zhang2024dynamic} or few-shot prompts \citep{han2024large}, and thus depend on hand-crafted design choices. Moreover, existing methods \citep{nam2024optimized,hollmann2023large} suffer from error propagation caused by simple feedback loops.

In contrast, our method does not rely on existing features, a predefined search space, or few-shot examples. Instead, it uses feature-importance scores to identify useful features rather than depending solely on performance evaluation. To reduce error propagation and ensure stable feature generation, we incorporate beam search, which produces more robust and consistent results.

\textbf{Pruning Search Space for CO Problems.} To improve the scalability of heuristics, several works have proposed to prune the search space of CO problems and focus only on the reduced ground set. These approaches fall into two categories: classical and learning-based algorithms.

\textbf{Classical Approaches.} \citet{zhou2017scaling} propose a pruning algorithm based on submodularity to reduce the cost of submodular maximization under size constraints. Similarly, \citet{pmlr-v258-nath25a} introduce a submodularity-based pruning algorithm that provides theoretical guarantees on both the fraction of the optimal value retained and the size of the resulting pruned ground set. The main limitation of these approaches is that they prune sequentially, which makes the process time-consuming. In contrast, our algorithm prunes in a single step. While this can produce a slightly larger pruned ground set in some cases, we demonstrate that our algorithm is orders of magnitude faster, particularly on large datasets.

\textbf{Learning-Based Approaches.} These algorithms can be divided into supervised learning methods and hybrid approaches that combine supervised and reinforcement learning. In the supervised setting, small instances are solved with exact or heuristic solvers, and the solutions are used to train a binary classifier that predicts whether a vertex or edge can be safely removed without compromising solution quality. These methods rely on hand-crafted, problem-specific features. For Maximum Clique Enumeration, \citet{lauri2019fine} apply logistic regression with hand-crafted features to prune vertices, while \citet{grassia2019learning,lauri2023learning} extends this line of work with a multi-stage pruning strategy in which successive classifiers, trained on graph-theoretic and statistical features, iteratively reduce the search space. Similarly, \citet{fitzpatrick2021learning,sun2021generalization} train classifiers to prune edges unlikely to belong to the optimal tour in Traveling Salesman Problem instance. 

The main limitation of these methods is their heavy reliance on problem-specific manual feature engineering and trial-and-error design, which restricts generalization. To address this, \citet{tian2024combhelper} proposes \comb, a GNN-based framework with knowledge distillation and a problem-specific boosting module designed to generalize across CO problems over graphs. While effective, it still depends on computationally expensive Fourier Feature Mapping and problem-specific boosting.


Hybrid methods that integrate supervised and reinforcement learning are generally more problem-agnostic. For instance, GCOMB \citep{manchanda2020gcomb} uses a probabilistic greedy algorithm to generate training data, followed by a linear regression model that applies a weighted-degree heuristic to prune the set. A Q-learning network is then trained on this reduced search space to sequentially predict the solution. A key limitation is that the initial pruning depends on a handcrafted degree heuristic, which does not generalize well across CO problems \citep{ireland2022lense}.

Similarly, \lense \citep{ireland2022lense} employs a two-stage learning process. First, a GNN is trained in a supervised manner to generate embeddings for subgraphs sampled from training graphs. Then, a policy network uses these embeddings to guide a local search procedure, moving from an initial subgraph to one more likely to contain the optimal solution. A major drawback, however, is its reliance on domain expertise to categorize subgraphs, set the number of samples per class, and tune critical hyperparameters such as subgraph size and embedding dimensionality separately for each dataset and problem.

As opposed to prior learning-based approaches, \llm makes the pruning process fully automated, requiring no domain knowledge, and can easily adapt to constraints.

% Hybrid methods that integrate supervised and reinforcement learning are generally more problem-agnostic. For instance, GCOMB \citep{manchanda2020gcomb} uses a probabilistic greedy algorithm for training data generation and then trains a linear regression model to generate a pruned set using a weighted degree heuristic. A Q-learning network is then trained on the pruned search space to sequentially predict the solution. A key limitation is that the initial pruning relies on a handcrafted degree heuristic, which does not generalize well across CO problems. Similarly, \lense \citep{ireland2022lense} employs a two-stage learning process. First, a GNN is trained in a supervised manner to generate embeddings for subgraphs sampled from training graphs. Then, a policy network uses these embeddings to guide a local search procedure, navigating from an initial subgraph to one more likely to contain the optimal solution. A major drawback is its reliance on domain expertise to categorize subgraphs, set the number of samples per class, and tune critical hyperparameters such as subgraph size and embedding dimensionality separately for each dataset and problem.

% \paragraph{Pruning Search Space for CO Problems.} To improve
% the scalability of heuristics, several works have proposed to
% prune the search space of CO problems and focus only on
% the reduced search space. These approaches primarily fall
% into two categories:  classical and learning-based algorithms.

% \textbf{Classical Approaches.} \cite{zhou2017scaling} propose a pruning algorithm based on the idea of submodularity to reduce the cost of submodular maximization under size constraints. Similarly, \cite{pmlr-v258-nath25a} introduce a pruning algorithm based on submodularity that provides theoretical guarantees on both the fraction of the optimal value retained and the size of the resulting pruned ground set. The main limitation of these approaches is that they prune the universe sequentially, which makes the process time-consuming. In contrast, our algorithm prunes the universe in a single step. While this produces a slightly larger pruned ground set and lesser objective value, we will show that our algorithms are faster by orders of magnitude.

% \textbf{Learning Based Approaches.} Learning-based approaches can be divided into two categories: supervised learning and a combination of supervised and reinforcement learning. In the supervised setting, small instances are solved with exact or heuristic solvers, and the optimal solutions are used to train a binary classifier that predicts whether a vertex or edge can be safely removed without compromising solution quality. These methods rely on hand-crafted features customized to each problem. For the Maximum Clique Enumeration problem, \citet{lauri2019fine} apply logistic regression with hand-crafted features to prune vertices, while \citet{grassia2019learning,lauri2023learning} extend this line of work with a multi-stage pruning strategy in which successive classifiers, trained on graph-theoretic and statistical features, are applied iteratively to progressively reduce the search space. \citet{fitzpatrick2021learning,sun2021generalization} train classifiers to prune edges unlikely to belong to the optimal tour in Traveling Salesman Problem instances. 
% The main limitation of these approaches is that they are heavily problem-specific, requiring substantial manual feature engineering and trial-and-error, which limits their generalization to other CO problems. 
% To address this limitation, \citet{tian2024combhelper} propose COMBHelper, a GNN-based framework with knowledge distillation and a problem-specific boosting module designed to generalize to multiple CO problems on graphs. While effective across various CO problems, it relies on computationally expensive Fourier Feature Mapping and a problem-specific boosting mechanism for strong performance. This problem-specific module leverages domain knowledge to enhance results—for example, in Minimum Vertex Cover, it penalizes the model more during training when it fails to identify high-degree nodes.

% Approaches that integrate supervised and reinforcement learning tend to be more problem-agnostic. For instance, GCOMB \citep{manchanda2020gcomb} uses a probabilistic greedy algorithm for training data generation and then trains a linear regression model to generate a pruned ground set using a weighted degree heuristic. A Q-learning network is then trained on the pruned search space to sequentially predict the solution. {However, a key limitation of this approach is that the initial pruning relies on a handcrafted degree heuristic, which does not generalize well to multiple CO problems .} Similarly, LeNSE \citep{ireland2022lense} employs a two-stage learning process: first, a GNN is trained using supervised learning to generate graph embeddings for subgraphs sampled from graphs in the training distribution. Then, a policy network is trained to use these embeddings to guide a local search procedure, navigating from an initial subgraph to a pruned subgraph that is more likely to contain the optimal solution. A major drawback of this method is its reliance on domain expertise to categorize subgraphs into different classes, determine the number of samples per class, and tune critical hyperparameters—such as the size of the vertex subset used to induce subgraphs and the dimensionality of vertex embeddings—individually for each dataset and problem to achieve strong performance.







